\documentclass[sigconf]{acmart}

\settopmatter{printacmref=false}
\setcopyright{none}
\renewcommand\footnotetextcopyrightpermission[1]{}

\usepackage{amsmath}
\usepackage{booktabs}

% --------------------------------------------------
% Paper information
% --------------------------------------------------

\title{NLP Assignment 1 Q1 Report}

\author{Muhammad Umer Siddiqui}
\email{ms08702@st.habib.pk}
\affiliation{%
  \institution{Dhanani School of Science and Engineering, Habib University}
  \city{Karachi}
  \country{Pakistan}
}

\author{Abdullah Shaikh}
\email{as09245@st.habib.pk}
\affiliation{%
  \institution{Dhanani School of Science and Engineering, Habib University}
  \city{Karachi}
  \country{Pakistan}
}
\author{Arqish} \email{az09714@st.habib.pk} \affiliation{% 
\institution{Dhanani School of Science and Engineering, Habib University} \city{Karachi} \country{Pakistan} }



\begin{document}

\begin{abstract}
This report presents implementations of Byte Pair Encoding (BPE) and
WordPiece tokenizers developed from scratch for a small multilingual
corpus containing English, Urdu, and Chinese. The report describes the
preprocessing pipeline, tokenization procedures, pair-selection criteria,
and the effect of vocabulary size on the learned subword vocabulary.
Experiments were conducted using vocabulary sizes ranging from 250 to
900. The results show that BPE began learning longer subword strings at
smaller vocabulary sizes, while WordPiece required a larger vocabulary
before producing substantial numbers of longer subword units. Based on
the observed results, a vocabulary size of 700 was selected as the final
configuration for the dataset.
\end{abstract}

\keywords{Natural Language Processing, Tokenization, BPE, WordPiece,
Multilingual NLP, Subword Tokenization}

\maketitle

% --------------------------------------------------
% Main content
% --------------------------------------------------

\section{Introduction}

The question provided a small multilingual corpus containing English,
Urdu, and Chinese text. The task was to implement Byte Pair Encoding
(BPE) and WordPiece tokenizers from scratch and experiment with their
hyperparameters on the given corpus. The assignment specifically
required consideration of the small corpus size, the three different
writing systems, code-switching, Unicode representation, and the
trade-off between vocabulary size and token fragmentation.

Our implementation used vocabulary size as its only explicit
hyperparameter. The algorithms were implemented in \texttt{q1.py}, and
the resulting vocabularies were stored in \texttt{bpe\_vocab.txt} and \\
\texttt{wordpiece\_vocab.txt}. We experimented with vocabulary sizes of
250, 300, 350, 450, 500, 600, 700, 800, and 900 to examine how the
learned subword units changed as the vocabulary grew.

\section{Preprocessing}

The corpus was first normalized using Unicode Normalization Form C
(NFC). This provides a canonical representation for characters that can
have multiple Unicode representations while preserving their linguistic
content. This was particularly important for a multilingual corpus
because the tokenizer operates directly on Unicode characters.

The text was then converted to lowercase. This primarily affected the
English portion of the corpus, where uppercase and lowercase versions
of the same letter would otherwise be treated as different characters.
Lowercasing did not produce a comparable change for the Urdu and
Chinese portions of the corpus.

Punctuation was separated from surrounding text rather than removed.
This allowed punctuation to remain available as tokens while preventing
it from being merged directly with adjacent words. Punctuation was
identified using Unicode character categories, making this step
independent of any particular language or writing system.

Finally, whitespace was normalized to single spaces. This provided a
consistent representation of word boundaries and allowed the
implementation to use whitespace-based splitting without being affected
by repeated spaces, tabs, or line breaks.

These preprocessing decisions were intended to normalize the corpus
without removing linguistic information. In particular, punctuation was
retained rather than discarded, and no stemming, lemmatization,
stopword removal, transliteration, or language-specific word
normalization was performed. Such transformations could alter the
frequency distribution of characters and subwords that the two
tokenizers use when deciding which pairs to merge.

\section{Tokenization Methods}

A parent class called \texttt{Tokenizer} was implemented to provide
functionality shared by both algorithms. It reads the corpus, performs
the preprocessing steps, and provides functionality for writing the
resulting vocabulary to a file. The BPE and WordPiece implementations
inherit from this class and implement their respective tokenization and
vocabulary-building procedures.

\subsection{WordPiece}

The \texttt{WordPieceTokenizer} class contains two main methods:
\texttt{tokenize} and \texttt{vocabulary\_builder}.

The \texttt{tokenize} method initially represents every word as a
sequence of individual characters. The first character of each word is
stored normally, while every subsequent character is prefixed with
\texttt{\#\#}. For example, the word \texttt{token} is initially
represented conceptually as

\[
[\texttt{t},\texttt{\#\#o},\texttt{\#\#k},\texttt{\#\#e},
\texttt{\#\#n}].
\]

The \texttt{vocabulary\_builder} method initializes the vocabulary using
these character-level tokens. It then repeatedly calculates the
frequency of each token and the frequency of every adjacent token pair.

Unlike BPE, WordPiece does not select a pair solely according to its
frequency. For two adjacent tokens \(x\) and \(y\), the implementation
uses the following score:

\[
\operatorname{Score}(x,y)
=
\frac{\operatorname{Freq}(x,y)}
{\operatorname{Freq}(x)\operatorname{Freq}(y)}.
\]

The pair with the highest score is selected for merging. The two tokens
are combined into a new token, with the \texttt{\#\#} marker removed
from the second token. All occurrences of the selected pair are then
replaced by the new token, which is added to the vocabulary.

This process continues until the target vocabulary size is reached or
no further adjacent pairs are available.

\subsection{Byte Pair Encoding}

The \texttt{BPETokenizer} class follows a similar iterative procedure
but uses a different pair-selection criterion.

The \texttt{tokenize} method splits each word into its individual
characters without adding WordPiece-style \texttt{\#\#} markers. For
example, the word \texttt{token} is initially represented as

\[
[\texttt{t},\texttt{o},\texttt{k},\texttt{e},\texttt{n}].
\]

The vocabulary is initialized using all unique characters in the
corpus. The algorithm then calculates the frequency of every adjacent
token pair. The pair with the highest raw frequency is selected:

\[
(x^*,y^*)
=
\underset{(x,y)}{\operatorname{arg\,max}}
\operatorname{Freq}(x,y).
\]

The selected pair is concatenated into a new token, and all occurrences
of the pair are replaced by the new token. The new token is added to
the vocabulary, and the process is repeated until the target vocabulary
size is reached.

\subsection{Comparison of Pair Selection}

The primary algorithmic difference between the two implementations is
therefore the criterion used to select pairs for merging. BPE selects
the pair with the highest raw frequency, whereas WordPiece normalizes
the pair frequency using the frequencies of the individual tokens.

This difference can cause the two algorithms to learn different
subword vocabularies even when they are given the same corpus and target
vocabulary size. BPE tends to favor pairs that occur frequently in the
corpus, while the WordPiece score can favor pairs that occur frequently
relative to the frequencies of their individual components.

\section{Experiments and Results}

The vocabulary size \(K\) was varied to examine how increasing the
available vocabulary affected the learned subwords. The tested values
were 250, 300, 350, 450, 500, 600, 700, 800, and 900.

\begin{table}[t]
\caption{Observed effect of vocabulary size on learned subwords}
\label{tab:vocab-results}
\centering
\begin{tabular}{c|p{2.6cm}|p{2.6cm}}
\toprule
$K$ & WordPiece & BPE \\
\midrule
250 & No token longer than one character was observed & Learned longer strings such as \texttt{word} and \texttt{tokeniz} \\
300 & Little observable change & More long strings, including \texttt{tokenization}, \texttt{process}, and \texttt{ing} \\
350 & Similar to $K=300$ & Continued improvement \\
450 & Noticeable improvement, particularly for Urdu & Continued improvement \\
500 & Little observable change & Little observable change \\
600 & Substantial number of English subword strings & Substantial number of complete words \\
700 & Further improvement & Further improvement \\
800 & Continued improvement & Continued improvement \\
900 & Continued improvement & Continued improvement \\
\bottomrule
\end{tabular}
\end{table}

At \(K=250\), WordPiece did not learn any token longer than one
character in the observed vocabulary. In contrast, BPE had already
learned longer strings such as \texttt{word} and \texttt{tokeniz}.
This indicates that the two pair-selection criteria produced
noticeably different vocabularies even at a relatively small
vocabulary size.

Increasing \(K\) to 300 produced little observable change in the
WordPiece vocabulary. BPE, however, learned additional longer strings,
including \texttt{tokenization} and \texttt{process}, as well as
suffixes such as \texttt{ing}. At \(K=350\), the observed results were
similar to those at \(K=300\).

At \(K=450\), WordPiece began producing substantially longer tokens,
with particularly noticeable changes in the Urdu portion of the
vocabulary. BPE also continued to produce longer subword units. The
increase to \(K=500\) produced comparatively little additional change
in either vocabulary.

At \(K=600\), WordPiece began learning a substantial number of English
subword strings. At the same vocabulary size, BPE was already learning
a substantial number of complete words. Increasing \(K\) to 700
continued to improve the learned vocabularies for both algorithms.
Further increases to 800 and 900 continued this trend for WordPiece.

\subsection{Multilingual Observations}

The multilingual nature of the corpus affected the vocabularies
produced by both algorithms. English words are composed from a
relatively small alphabet and contain recurring character sequences,
allowing frequent sequences to be merged into longer subwords. This is
visible in the English examples learned by BPE at relatively small
vocabulary sizes.

The Urdu portion of the corpus showed a more noticeable improvement
for WordPiece at larger vocabulary sizes. This demonstrates that the
choice of vocabulary size can affect the representation of different
languages differently within a shared multilingual vocabulary.

Chinese differs from English and Urdu because its writing system does
not use alphabetic word formation in the same way. Consequently, the
character-level initialization provides a particularly direct
representation of Chinese text. The learned Chinese vocabulary should
therefore be interpreted in terms of character sequences rather than
assuming that every learned token corresponds to an English-style
word.

\subsection{BPE and WordPiece Comparison}

The experiments demonstrate a difference in how the two algorithms use
additional vocabulary capacity. BPE began producing longer strings at
smaller vocabulary sizes in our experiments, while WordPiece showed a
more noticeable increase in longer subword strings at larger values of
\(K\).

The difference follows from their pair-selection criteria. BPE directly
maximizes pair frequency, so frequently occurring sequences are merged
first. WordPiece instead considers pair frequency relative to the
frequencies of the individual tokens. As a result, a pair that is
frequent in isolation is not necessarily selected by WordPiece at the
same stage at which it would be selected by BPE.

\section{Discussion}

The experiments show that vocabulary size has a direct effect on the
degree of token fragmentation. A small vocabulary limits the number of
merges that can be performed and therefore leaves more of the corpus
represented by individual characters or short subwords. Increasing the
vocabulary size allows additional merges and consequently produces
longer learned units.

The results also show that increasing \(K\) does not necessarily
produce the same degree of improvement at every step. For example,
increasing \(K\) from 250 to 300 produced a more noticeable change for
BPE than for WordPiece, while larger vocabulary sizes produced more
substantial changes in WordPiece.

There is therefore a trade-off in selecting \(K\). A very small
vocabulary results in greater token fragmentation, whereas a very large
vocabulary can allocate vocabulary entries to increasingly specific
strings. Since the assignment corpus is small and multilingual, the
selected value should provide useful subword units across the different
languages rather than simply maximizing the number of learned strings.

Based on the observed results, \(K=700\) was selected as the final
configuration. At this value, WordPiece had begun learning a
substantial number of longer subword strings, while BPE was producing
many complete words and longer subword units. Values above 700
continued the observed trend, but the additional vocabulary capacity
was not considered necessary for the purposes of this small corpus.

\section{Conclusion}

This report implemented BPE and WordPiece from scratch and examined the
effect of vocabulary size on their learned subword vocabularies. The
two algorithms differed primarily in their pair-selection criteria:
BPE selected the most frequent adjacent pair, while WordPiece used a
normalized score based on pair and individual-token frequencies.

Experiments across vocabulary sizes from 250 to 900 showed that BPE
began learning longer strings at smaller vocabulary sizes, while
WordPiece showed more substantial improvements at larger vocabulary
sizes. The experiments also demonstrated that vocabulary size affects
different languages differently within a shared multilingual
vocabulary.

Based on the observed results, a vocabulary size of 700 was selected
for the final configuration because it provided a reasonable balance
between vocabulary size and the quality of the learned subword units
for the given corpus.

\end{document}